\documentclass[11pt]{article}

\usepackage[margin=1in]{geometry}
\usepackage[T1]{fontenc}
\usepackage[utf8]{inputenc}
\usepackage{lmodern}
\usepackage{booktabs}
\usepackage{longtable}
\usepackage{array}
\usepackage{siunitx}
\usepackage{hyperref}
\usepackage{caption}

\sisetup{round-mode=places,round-precision=4}

\title{Final Study Report\\Curriculum Learning with Difficulty-Aware Negative Sampling for GNN Link Prediction}
\author{IAI Research Project}
\date{5 April 2026}

\begin{document}

\maketitle

\begin{abstract}
This report presents the final consolidated results for a graph link-prediction study comparing standard random-negative training with curriculum learning based on structural hard-negative sampling. The pipeline includes GCN and GAT models, three difficulty heuristics (CN, AA, RA), adaptive curriculum scheduling, and HeaRT evaluation. Across full multi-seed runs (baseline: 60, curriculum: 180, ablation: 80), findings are mixed and dataset-dependent: curriculum variants show small gains on PubMed in HeaRT MRR but underperform baseline on Cora and Citeseer under current presets. The strongest degradations occur for Cora with CN-focused curriculum, while baseline or baseline-like settings remain strongest in the Cora+GCN ablation matrix.
\end{abstract}

\section{Introduction}

Link prediction with graph neural networks (GNNs) is sensitive to negative-sampling strategy. Uniform random negatives are often easy and can inflate apparent ranking quality. This project evaluates whether curriculum learning over difficulty-aware negatives improves robustness under hard-negative evaluation.

The main research question is:

\begin{quote}
Can adaptive curriculum learning over structurally hard negatives improve HeaRT performance relative to random-negative baseline training?
\end{quote}

\section{Method Summary}

\subsection{Models and Data}

\begin{itemize}
\item Models: GCN and GAT.
\item Datasets: Cora, Citeseer, PubMed.
\item Split strategy: link-prediction split generated from the graph with fixed seeds.
\end{itemize}

\subsection{Difficulty-Aware Negatives and Curriculum}

Candidate negative edges are precomputed and scored by one structural heuristic:
\begin{itemize}
\item CN (common neighbors)
\item AA (Adamic-Adar)
\item RA (resource allocation)
\end{itemize}

Candidates are bucketed into easy/medium/hard partitions. The default adaptive curriculum uses four phases:
\begin{center}
easy\_only $\rightarrow$ easy\_medium $\rightarrow$ mixed $\rightarrow$ hard\_focus.
\end{center}

\subsection{Evaluation}

Two metric families are reported:
\begin{itemize}
\item Standard: AUC, AP, MRR, Hits@K
\item HeaRT: hard-negative ranking metrics (HeaRT MRR, HeaRT Hits@K)
\end{itemize}

\section{Experimental Matrix}

\begin{itemize}
\item Baseline runs: $3$ datasets $\times$ $2$ models $\times$ $10$ seeds $= 60$.
\item Curriculum runs: $3$ datasets $\times$ $2$ models $\times$ $3$ heuristics $\times$ $10$ seeds $= 180$.
\item Ablation runs (Cora+GCN): $8$ conditions $\times$ $10$ seeds $= 80$.
\end{itemize}

All summaries are taken from:
\begin{itemize}
\item \texttt{results/summaries/baseline\_summary.csv}
\item \texttt{results/summaries/curriculum\_summary.csv}
\item \texttt{results/summaries/ablation\_summary.csv}
\item \texttt{results/summaries/full\_significance\_table.csv}
\item \texttt{results/summaries/significance\_tests.csv}
\end{itemize}

\section{Main Results}

\subsection{Baseline vs Best Curriculum (HeaRT MRR)}

\begin{table}[h]
\centering
\caption{Baseline vs best curriculum heuristic by dataset/model using 10-seed mean HeaRT MRR.}
\begin{tabular}{llcccl}
\toprule
Dataset & Model & Baseline & Best Curriculum & Best Heuristic & Relative Change \\
\midrule
Cora & GCN & 0.5055 & 0.4124 & AA & -18.43\% \\
Cora & GAT & 0.5226 & 0.5140 & RA & -1.65\% \\
Citeseer & GCN & 0.4761 & 0.4593 & CN & -3.53\% \\
Citeseer & GAT & 0.5781 & 0.5435 & RA & -6.00\% \\
PubMed & GCN & 0.5420 & 0.5480 & AA & +1.12\% \\
PubMed & GAT & 0.4752 & 0.4769 & RA & +0.37\% \\
\bottomrule
\end{tabular}
\end{table}

\subsection{Heuristic-Level HeaRT MRR (Curriculum)}

\begin{table}[h]
\centering
\caption{Curriculum HeaRT MRR means by heuristic.}
\begin{tabular}{llccc}
\toprule
Dataset & Model & CN & AA & RA \\
\midrule
Cora & GCN & 0.3854 & 0.4124 & 0.4104 \\
Cora & GAT & 0.3725 & 0.5061 & 0.5140 \\
Citeseer & GCN & 0.4593 & 0.4585 & 0.4577 \\
Citeseer & GAT & 0.5297 & 0.5328 & 0.5435 \\
PubMed & GCN & 0.5475 & 0.5480 & 0.5469 \\
PubMed & GAT & 0.4760 & 0.4754 & 0.4769 \\
\bottomrule
\end{tabular}
\end{table}

\section{Ablation Study (Cora + GCN)}

\begin{table}[h]
\centering
\caption{Ablation ranking by HeaRT MRR (higher is better).}
\begin{tabular}{lcc}
\toprule
Condition & HeaRT MRR & Relative to \texttt{abl-1} \\
\midrule
\texttt{abl-1} (baseline control) & 0.5051 & 0.00\% \\
\texttt{abl-10} (low thresholds) & 0.3880 & -23.17\% \\
\texttt{abl-6} (extra intermediate phase) & 0.3816 & -24.46\% \\
\texttt{abl-5} (skip mixed) & 0.3811 & -24.54\% \\
\texttt{abl-3} (static easy/hard) & 0.3664 & -27.47\% \\
\texttt{abl-11} (high thresholds) & 0.3634 & -28.06\% \\
\texttt{abl-2} (hard from start) & 0.3498 & -30.75\% \\
\texttt{abl-4} (fixed-phase schedule) & 0.3419 & -32.31\% \\
\bottomrule
\end{tabular}
\end{table}

\section{Statistical Results}

From \texttt{full\_significance\_table.csv}:

\begin{itemize}
\item Total comparison rows: 69
\item HeaRT MRR rows: 23
\item Significant HeaRT MRR rows ($p<0.05$): 4
\end{itemize}

The 4 significant HeaRT MRR effects are:

\begin{table}[h]
\centering
\caption{Significant HeaRT MRR curriculum-vs-baseline comparisons (from full significance table).}
\begin{tabular}{llcccc}
\toprule
Dataset & Model & Heuristic & Improvement (\%) & p-value & Direction \\
\midrule
Citeseer & GAT & CN & -8.38 & 0.0475 & Worse \\
Cora & GAT & CN & -28.72 & $3.44\times10^{-6}$ & Worse \\
Cora & GCN & AA & -18.43 & 0.00135 & Worse \\
Cora & GCN & RA & -18.81 & 0.00127 & Worse \\
\bottomrule
\end{tabular}
\end{table}

\section{Discussion}

\subsection{What Worked}
\begin{itemize}
\item Curriculum variants were competitive on PubMed, with small positive HeaRT MRR gains for GCN and GAT in best-heuristic settings.
\item The codebase successfully supports full multi-seed reproducible baselines, curricula, ablations, significance analysis, and figure generation.
\end{itemize}

\subsection{What Did Not Work Under Current Presets}
\begin{itemize}
\item On Cora and Citeseer, best curriculum settings underperform baseline in HeaRT MRR.
\item The strongest degradation appears in CN-focused curricula on Cora, especially for GAT.
\item Cora+GCN ablations indicate baseline-like control remains strongest among tested variants.
\end{itemize}

\subsection{Interpretation}
The evidence suggests strong dataset dependence. Hard-negative exposure schedules that appear reasonable globally are not universally beneficial. Threshold design, phase composition, and heuristic choice likely require dataset-specific tuning to avoid over-regularizing or destabilizing ranking behavior.

\section{Limitations}

\begin{itemize}
\item Curriculum presets were evaluated without extensive per-dataset hyperparameter optimization.
\item Statistical testing is reported from existing scripts; multiple-comparison correction is not currently integrated in the pipeline.
\item Current conclusions are bounded to tested datasets (Cora, Citeseer, PubMed) and model family (GCN/GAT).
\end{itemize}

\section{Conclusion}

The final study does not support a universal claim that the current curriculum presets improve hard-negative link prediction performance. Results are mixed: slight improvements on PubMed and consistent degradations in key Cora/Citeseer settings. The framework itself is complete and reproducible, and provides a solid basis for next-stage work focused on adaptive threshold tuning, alternative schedules, and broader model/data coverage.

\section*{Reproducibility}

Primary commands:
\begin{verbatim}
./scripts/run_all_experiments.sh baseline
./scripts/run_all_experiments.sh curriculum
./scripts/run_all_experiments.sh ablation
./scripts/run_all_experiments.sh aggregate
python experiments/statistical_analysis.py
python experiments/generate_figures.py
\end{verbatim}

Compile this report with:
\begin{verbatim}
cd docs
pdflatex final_study_report.tex
\end{verbatim}

\section*{Reference}

\begin{thebibliography}{1}
\bibitem{heart2023}
Benjamin Chamberlain et al.
\newblock HeaRT: Hard Evaluation for Graph Representation Learning.
\newblock Advances in Neural Information Processing Systems, 2023.
\end{thebibliography}

\end{document}